These are Freeks personal critical notes when going over the text. 

Lemma 2.1.3. What does closed mean? One can also take this as definition of closed and make a remark similar to 2.1.2. 

Similar remark on 2.2.3. 

General citation style (Dfn 2.5 vs Definition, parencite vs cite). 

Do we mention the axioms of synthetic-stone-duality before we start?
Maybe general statement at the start of this section should be that we write down only results that we use in the rest of the paper, and might have been rephrased or slightly expanded in some direction. In the sections themselves, we refer to synthetic-stone-duality for the proofs, and we don't copy all statement/axioms unless we use them in the rest of the paper. 
TODO when done reading: check that indeed all the definition/lemmata in these sections are used in the rest of the paper. 

Lemma 2.1.4.iii is it somehow clear that odisc types are closed under sequential colimits? I feel it's equally complicated. I tought the argument was dependent choice to pick compatible representations of finite sets and then do a diagonal argument, which is quite complicated for this part, but might be already in the original somewhere? No it's not, but we should refer to the lemma that sequential colimits of Odisc is ODisc. Ohno, it is obvious that sequential colimits of countably presented boolean algebras are still countably presented boolean algebras. 
Conclusion, there are other proofs available, this one works and highlights duality. One other proof could indeed work with carefully chosen presentations and dependent choice to make a diagonal argument. and then some limit argument. 

Lemma 2.1.7 I removed some extra enters. That's a style discussion: how much space do we take?

Lemma 2.2.5.v what does that mean? 
Why is 2.2.5.iv obvious? Is it obvious that if I quotient by an open equivalence relation, and then by a new open equivalence relation, then I could just as well quotient by a composite relation? I have no proof that open relations compose neatly (equivalence relations I think should)

TODO ask Reid how we should cite Barton-Commelin condensed type theory. 

2.2.6 the "actually holds for any finitary algebraic notion" remark in the proof suggests this Lemma and proof should be a remark instead. 
\begin{remark}
  In Corollary 2.14 of \cite{synthetic-stone-duality}, it is shown that a Boolean algebra is countably presented if and only if it is overtly discrete. Such a result holds if we replace "Boolean algebra" by any finiteary algebraic notion, such as group or Abelian group. 
\end{remark}
